\documentclass[12pt, a4paper]{article}

\usepackage{amsmath, amssymb, amsfonts}
\usepackage{float}
\usepackage{graphicx}
\usepackage{listings}
\usepackage{xcolor}
\usepackage[linesnumbered, ruled]{algorithm2e}
\usepackage{subcaption}
\usepackage[hidelinks]{hyperref}
% Each algorithm is written twice: the algorithm2e block (what the PDF typesets)
% and a plain-text Verbatim block right after it (what the web shows — pandoc
% renders Verbatim as a copyable <pre> code block). \excludecomment makes pdflatex
% SKIP the Verbatim blocks, so the PDF shows only the algorithm2e version; pandoc
% conversely drops each algorithm2e to an (empty, CSS-hidden) div and keeps the
% Verbatim. Because pdflatex never sees the Verbatim, its Unicode maths glyphs
% (⟨⟩ ← τ ∈ ∞ …) need no inputenc/literate setup.
\usepackage{comment}
\excludecomment{Verbatim}

\definecolor{codegreen}{rgb}{0,0.6,0}
\definecolor{codegray}{rgb}{0.5,0.5,0.5}
\definecolor{codepurple}{rgb}{0.58,0,0.82}
\hypersetup{linktoc=all}

\lstset{
    language=Python,
    basicstyle=\ttfamily\small,
    keywordstyle=\color{blue},
    commentstyle=\color{codegreen},
    stringstyle=\color{codepurple},
    breaklines=true,
    numbers=left,
    numberstyle=\tiny\color{codegray},
    frame=single,
    showstringspaces=false
}

\title{AI Masterclass Test 1}
\author{}
\date{}


\begin{document}

\section{Minimax Search}
\begin{algorithm}[H]
\caption{Minimax search}
\SetKwFunction{Fback}{backpropagate}
\SetKwProg{Fn}{function}{}{end function}

$done \leftarrow leaves((S, E))$\;
\While{$done \neq S$}{
    $next \leftarrow \{s' \in S \setminus done \mid children(s', (S, E)) \subseteq done\}$\;
    \For{$s \in next$}{
        $C \leftarrow children(s, (S, E))$\;
        \If{$owner(s) = 1$}{
            $O \leftarrow \arg \max_{s' \in C} u(s')$ \tcp*[r]{maximise}
        }
        \Else{
            $O \leftarrow \arg \min_{s' \in C} u(s')$ \tcp*[r]{minimise}
        }
        $s'' \leftarrow$ any element of $O$\;
        $u(s) \leftarrow u(s'')$ \tcp*[r]{back score up}
        $done \leftarrow done \cup \{s\}$\;
    }
}
\end{algorithm}
\begin{Verbatim}
# Minimax search
done ← leaves((S, E))
while done ≠ S:
    next ← { s' ∈ S ∖ done | children(s', (S, E)) ⊆ done }
    for s ∈ next:
        C ← children(s, (S, E))
        if owner(s) = 1:
            O ← arg max_{s' ∈ C} u(s')        # maximise
        else:
            O ← arg min_{s' ∈ C} u(s')        # minimise
        s'' ← any element of O
        u(s) ← u(s'')                         # back score up
        done ← done ∪ {s}
\end{Verbatim}

\begin{itemize}
    \item Here, we are concerned with zero sum games, which are a subset of extended form games.
    \item For a two player zero sum game (where one player's win is the other player's loss)
    \[G = (N = \{1, 2\},\: (S, E, s_0),\: owner,\: a,\: u)\]
    \begin{itemize}
        \item The components are almost the same as extended form games, except:
        \item $N = \{1,2\}$ is set of two players
        \item $u: leaves((S, E))\rightarrow \{1,-1\}$ gives the scores for player 1.
        The score of player 2 is the negation of score of player 1. 
    \end{itemize}
    \item Minimax is a variation of backward induction, optimised for zero sum games.
\end{itemize}

\subsection{Depth-limited minimax with heuristics} In practice, real game trees are too large for full backward induction.
We want to use depth-limited search. We then have to replace the exact utilities at new leaves with some heuristic evaluations.
For example, in chess it might be $h(s) = 9Q + 5R + 3B$.

\subsection{Alpha-beta pruning}
An optimisation for minimax. 
\begin{figure}[H]
    \centering
    \includegraphics[width=0.8\textwidth]{images/alpha-beta-pruning.png} 
    % \caption{$h(s)$ = number of misplaced tiles or sum of Manatten distances of each tile to their goal position}
    \label{fig:my_handwritten_notes}
\end{figure}
\begin{itemize}
    \item Keep track of $\alpha$, the best (highest) value that Maximiser has found so far.
    \item $\beta$ is the lowest value that Maximiser has found so far.
    \item If $\beta \leq \alpha$, stop searching a branch. There is no point as Maximiser cannot get a higher value $\alpha$.
\end{itemize}

\end{document}
